\section{AIME I 1989}
        \begin{enumerate}
        	\item (\textit{AIME I 1989}) Compute $\sqrt{(31)(30)(29)(28)+1}$.


 
 \textbf{Answer: }869



\textbf{Solution 1 (Symmetry)}

Note that the four numbers to multiply are symmetric with the center at $29.5$. Multiply the symmetric pairs to get $31\cdot 28=868$ and $30\cdot 29=870$.$\sqrt{868\cdot 870 + 1} = \sqrt{(869-1)(869+1) + 1} = \sqrt{869^2 - 1^2 + 1} = \sqrt{869^2} = \boxed{869}$.

 





\textbf{Solution 2 (Symmetry)}

Notice that $(a+1)^2 = a \cdot (a+2) +1$. Then we can notice that $30 \cdot 29 =870$ and that $31 \cdot 28 = 868$. Therefore, $\sqrt{(31)(30)(29)(28) +1} = \sqrt{(870)(868) +1} = \sqrt{(868 +1)^2} = \boxed{869}$. This is because we have that $a=868$ as per the equation $(a+1)^2 = a \cdot (a+2) +1$.

 





\textbf{Solution 3 (Symmetry with Generalization)}

More generally, we can prove that one more than the product of four consecutive integers must be a perfect square:
 \begin{align*} (a+3)(a+2)(a+1)(a)+1 &= [(a+3)(a)][(a+2)(a+1)]+1 \\ &= [a^2+3a][a^2+3a+2]+1 \\ &= [a^2+3a]^2+2[a^2+3a]+1 \\ &= [a^2+3a+1]^2. \end{align*}
At $a=28,$ we have 
 \[\sqrt{(a+3)(a+2)(a+1)(a)+1}=a^2+3a+1=\boxed{869}.\]~Novus677 ~MRENTHUSIASM

 





\textbf{Solution 4 (Symmetry with Generalization)}

Similar to Solution 1 above, call the consecutive integers $\left(n-\frac{3}{2}\right), \left(n-\frac{1}{2}\right), \left(n+\frac{1}{2}\right), \left(n+\frac{3}{2}\right)$ to make use of symmetry.  Note that $n$ itself is not an integer - in this case, $n = 29.5$.  The expression becomes $\sqrt{\left(n-\frac{3}{2}\right)\left(n + \frac{3}{2}\right)\left(n - \frac{1}{2}\right)\left(n + \frac{1}{2}\right) + 1}$.  Distributing each pair of difference of squares first, and then distributing the two resulting quadratics and adding the constant, gives $\sqrt{n^4 - \frac{5}{2}n^2 + \frac{25}{16}}$.  The inside is a perfect square trinomial, since $b^2 = 4ac$.  It's equal to $\sqrt{\left(n^2 - \frac{5}{4}\right)^2}$, which simplifies to $n^2 - \frac{5}{4}$.  You can plug in the value of $n$ from there, or further simplify to $\left(n - \frac{1}{2}\right)\left(n + \frac{1}{2}\right) - 1$, which is easier to compute.  Either way, plugging in $n=29.5$ gives $\boxed{869}$.

 





\textbf{Solution 5 (Prime Factorization)}

We have $(31)(30)(29)(28)+1=755161.$ Since the alternating sum of the digits $7-5+5-1+6-1=11$ is divisible by $11,$ we conclude that $755161$ is divisible by $11.$

 We evaluate the original expression by prime factorization:
 \begin{align*} \sqrt{(31)(30)(29)(28)+1}&=\sqrt{755161} \\ &=\sqrt{11\cdot68651} \\ &=\sqrt{11^2\cdot6241} \\ &=\sqrt{11^2\cdot79^2} \\ &=11\cdot79 \\ &=\boxed{869}. \end{align*}
~Vrjmath ~MRENTHUSIASM

 





\textbf{Solution 6 (Observation)}

The last digit under the radical is $1$, so the square root must either end in $1$ or $9$, since $x^2  = 1\pmod {10}$ means $x = \pm 1$.  Additionally, the number must be near $29 \cdot 30 = 870$, narrowing the reasonable choices to $869$ and $871$.

 Continuing the logic, the next-to-last digit under the radical is the same as the last digit of $28 \cdot 29 \cdot 3 \cdot 31$, which is $6$.  Quick computation shows that $869^2$ ends in $61$, while $871^2$ ends in $41$.  Thus, the answer is $\boxed{869}$.

 





\textbf{Solution 7 (Brute force)}

$31 \times 30 \times 29 \times 28$ is $755160$. $755160 + 1=755161$, and $\sqrt{755161}=869$. So the answer is $\boxed{869}$.

 




	\item (\textit{AIME I 1989}) Ten points are marked on a circle. How many distinct convex polygons of three or more sides can be drawn using some (or all) of the ten points as vertices?


 
 \textbf{Answer: }968



\textbf{Solution}

Any \href{https://artofproblemsolving.com/wiki/index.php?title=Subset}{subset} of the ten points with three or more members can be made into exactly one such polygon.  Thus, we need to count the number of such subsets.  There are $2^{10} = 1024$ total subsets of a ten-member \href{https://artofproblemsolving.com/wiki/index.php?title=Set}{set}, but of these ${10 \choose 0} = 1$ have 0 members, ${10 \choose 1} = 10$ have 1 member and ${10 \choose 2} = 45$ have 2 members.  Thus the answer is $1024 - 1 - 10 - 45 = \boxed{968}$.

 Note ${n \choose 0}+{n \choose 1} + {n \choose 2} + \dots + {n \choose n}$ is equivalent to $2^n$

 




	\item (\textit{AIME I 1989}) Suppose $n_{}^{}$ is a positive integer and $d_{}^{}$ is a single digit in base 10. Find $n_{}^{}$ if


 $\frac{n}{810}=0.d25d25d25\ldots$
 \textbf{Answer: }750



\textbf{Solution 1}

We can express $0.\overline{d25}$ as $\frac{100d+25}{999}$. We set up the given equation and isolate $n:$ \begin{align*} \frac{100d+25}{999} &= \frac{n}{810}, \\ \frac{100d+25}{111} &= \frac{n}{90}, \\ 9000d + 2250 &= 111n. \end{align*}
 We then set up the following modular congruence to solve for $d:$ \begin{align*}9000d + 2250 &\equiv 0 \pmod {111}, \\ 9d + 30 &\equiv 0 \pmod {111}, \\ 3d + 10 &\equiv 0 \pmod {37}, \\ 36d + 120 &\equiv 0 \pmod {37}, \\ -d &\equiv -9 \pmod {37}, \\ d &\equiv 9 \pmod {37}.\end{align*}
 Since $d$ is a digit, it must be equal to $9$ based on our above constraint. When $d=9, n = \boxed{750}.$

 





\textbf{Solution 2}

Repeating decimals represent \href{https://artofproblemsolving.com/wiki/index.php?title=Rational\_number}{rational numbers}.  To figure out which rational number, we sum an \href{https://artofproblemsolving.com/wiki/index.php?title=Infinite}{infinite} \href{https://artofproblemsolving.com/wiki/index.php?title=Geometric\_series}{geometric series}, $0.d25d25d25\ldots = \sum_{n = 1}^\infty \frac{d25}{1000^n} = \frac{100d + 25}{999}$.  Thus $\frac{n}{810} = \frac{100d + 25}{999}$ so $n = 30\frac{100d + 25}{37} =750\frac{4d + 1}{37}$.  Since 750 and 37 are \href{https://artofproblemsolving.com/wiki/index.php?title=Relatively\_prime}{relatively prime}, $4d + 1$ must be \href{https://artofproblemsolving.com/wiki/index.php?title=Divisible}{divisible} by 37, and the only digit for which this is possible is $d = 9$.  Thus $4d + 1 = 37$ and $n = \boxed{750}$.

 
(Note: Any repeating sequence of $n$ digits that looks like $0.a_1a_2a_3...a_{n-1}a_na_1a_2...$ can be written as $\frac{a_1a_2...a_n}{10^n-1}$, where $a_1a_2...a_n$ represents an $n$ digit number.)

 





\textbf{Solution 3}

To get rid of repeating decimals, we multiply the equation by 1000. We get $\frac{1000n}{810} = d25.d25d25...$ We subtract the original equation from the second to get $\frac{999n}{810}=d25$ We simplify to $\frac{37n}{30} = d25$ Since $\frac{37n}{30}$ is an integer, $n=(30)(5)(2k+1)$ because $37$ is relatively prime to $30$, and d25 is divisible by $5$ but not $10$. The only odd number that yields a single digit $d$ and 25 at the end of the three digit number is $k=2$, so the answer is $\boxed{750}$.

 





\textbf{Solution 4}

Similar to Solution 2, we start off by writing that $\frac{1000n}{810} = d25.d25d25 \dots$ .Then we subtract this from the original equation to get:

 
 \[\frac{999n}{810} =d25 \Longrightarrow \frac{37n}{30} = d25 \Longrightarrow 37n = d25 \cdot 30\]

 Since n is an integer, we have that $37 \mid d25 \cdot 30$.

 Since $37$ is prime, we can apply \href{https://artofproblemsolving.com/wiki/index.php?title=Euclid\%27s\_Lemma}{Euclid's Lemma} to realize that $37 \mid d25$, since $37 \nmid 30$. Then we can expand $d25$ as $25 \cdot (4d +1)$. Since $37 \nmid 25$, by Euclid, we can arrive at $37 \mid 4d+1 \Longrightarrow d=9$. From this we know that $n= 25 \cdot 30 = \boxed{750}$. (This is true because $37n = 925 \cdot 30 \rightarrow n= 25 \cdot 30 = 750$)

 





\textbf{Solution 5}

Write out these equations:

 
$\frac{n}{810} = \frac{d25}{999}$

 
$\frac{n}{30} = \frac{d25}{37}$

 
$37n = 30(d25)$

 
Thus $n$ divides 25 and 30. The only solution for this under 1000 is $\boxed{750}$.

 
-jackshi2006

 Note: We know $n < 810$ since $\frac{n}{810} < 1,$ so it suffices to check for numbers under $810$ not $1000.$

 




	\item (\textit{AIME I 1989}) If $a<b<c<d<e^{}_{}$ are consecutive positive integers such that $b+c+d$ is a perfect square and $a+b+c+d+e^{}_{}$ is a perfect cube, what is the smallest possible value of $c^{}_{}$?


 
 \textbf{Answer: }675



\textbf{Solution}

Since the middle term of an \href{https://artofproblemsolving.com/wiki/index.php?title=Arithmetic\_progression}{arithmetic progression} with an odd number of terms is the average of the series, we know $b + c + d = 3c$ and $a + b + c + d + e = 5c$. Thus, $c$ must be in the form of $3 \cdot x^2$ based upon the first part and in the form of $5^2 \cdot y^3$ based upon the second part, with $x$ and $y$ denoting an \href{https://artofproblemsolving.com/wiki/index.php?title=Integer}{integers}. $c$ is minimized if it’s \href{https://artofproblemsolving.com/wiki/index.php?title=Prime\_factorization}{prime factorization} contains only $3,5$, and since there is a cubed term in $5^2 \cdot y^3$, $3^3$ must be a factor of $c$. $3^35^2 = \boxed{675}$, which works as the solution.

 





\textbf{Solution 2}

Let $b$, $c$, $d$, and $e$ equal $a+1$, $a+2$, $a+3$, and $a+4$, respectively. Call the square and cube $k^2$ and $m^3$, where both k and m are integers. Then:

 $5a + 10 = m^3$

 Now we know $m^3$ is a multiple of 125 and $m$ is a multiple of 5. The lower $m$ is, the lower the value of $c$ will be. Start from 5 and add 5 each time.

 $m = 5$ gives no solution for k

 $m = 10$ gives no solution for k

 $m = 15$ gives a solution for k.

 
$10 + 5a = 15^3$

 
$2 + a = 675$

 
$c = \boxed{675}$

 
-jackshi2006

 


 





\textbf{Solution 3}

Let the numbers be $a,a+1,a+2,a+3,a+4.$ When then know $3a+6$ is a perfect cube and $5a+10$ is perfect cube. Since $5a+10$ is divisible by $5$ we know that $5a+10 = (5k)^3$ since otherwise we get a contradiction. This means $a = 25k^3 - 2$ in which plugging into the other expression we know $3(25k^3 - 2) + 6 = 75k^3$ is a perfect square. We know $75 = 5^2 \cdot 3$ so we let $k = 3$ to obtain the perfect square. This means that $c = a+2 = (25 \cdot 27 - 2)+2 = 25 \cdot 27 = \boxed{675}.$

 





\textbf{Solution 4}

(This is literally a combination of 1 and 3) 

 Since $a$, $b$, $c$, $d$, and $e$ are consecutive, $a = c-2$, $b = c-1$, $c=c$, $d = c+1$, and $e = c+2$. 

 Because $b+c+d = 3c$ is a perfect square, and $a+b+c+d+e = 5c$ is a perfect cube, we can express $c$ as $c = 3^{n} \cdot 5^{k}$. 

 Now, by the problem's given information, 

 $k \equiv 0   \text{(mod 2)}$

 $n \equiv 0   \text{(mod 3)}$ 

 and because ALL exponents have to be cubes/squares, 

 $k \equiv 2   \text{(mod 3)}$

 $n \equiv 1   \text{(mod 2)}$ 

 Therefore, 

 $k = 2$

 $n = 3$ 

 $c = 3^3 \cdot 5^2 = \boxed{675}$.

 




	\item (\textit{AIME I 1989}) When a certain biased coin is flipped five times, the probability of getting heads exactly once is not equal to $0$ and is the same as that of getting heads exactly twice. Let $\frac ij^{}_{}$, in lowest terms, be the probability that the coin comes up heads in exactly $3$ out of $5$ flips. Find $i+j^{}_{}$.


 
 \textbf{Answer: }283



\textbf{Solution 1}

Denote the probability of getting a heads in one flip of the biased coin as $h$. Based upon the problem, note that ${5\choose1}(h)^1(1-h)^4 = {5\choose2}(h)^2(1-h)^3$. After canceling out terms, we get $1 - h = 2h$, so $h = \frac{1}{3}$. The answer we are looking for is ${5\choose3}(h)^3(1-h)^2 = 10\left(\frac{1}{3}\right)^3\left(\frac{2}{3}\right)^2 = \frac{40}{243}$, so $i+j=40+243=\boxed{283}$.

 


 





\textbf{Solution 2}

Denote the probability of getting a heads in one flip of the biased coins as $h$ and the probability of getting a tails as $t$. Based upon the problem, note that ${5\choose1}(h)^1(t)^4 = {5\choose2}(h)^2(t)^3$. After cancelling out terms, we end up with $t = 2h$. To find the probability getting $3$ heads, we need to find ${5\choose3}\dfrac{(h)^3(t)^2}{(h + t)^5} =10\cdot\dfrac{(h)^3(2h)^2}{(h + 2h)^5}$ (recall that $h$ cannot be $0$). The result after simplifying is $\frac{40}{243}$, so $i + j = 40 + 243 = \boxed{283}$.

 




	\item (\textit{AIME I 1989}) Two skaters, Allie and Billie, are at points $A$ and $B$, respectively, on a flat, frozen lake. The distance between $A$ and $B$ is $100$ meters. Allie leaves $A$ and skates at a speed of $8$ meters per second on a straight line that makes a $60^\circ$ angle with $AB$. At the same time Allie leaves $A$, Billie leaves $B$ at a speed of $7$ meters per second and follows the straight path that produces the earliest possible meeting of the two skaters, given their speeds. How many meters does Allie skate before meeting Billie?


 
                    \begin{figure}[H]
                        \centering
                        \includegraphics[width=0.5\textwidth]{Problem/AIME_I_1989_6_Problem_Figure_1.png}
                    \end{figure}
                    
 \textbf{Answer: }160



\textbf{Solution}

Label the point of \href{https://artofproblemsolving.com/wiki/index.php?title=Intersection}{intersection} as $C$. Since $d = rt$, $AC = 8t$ and $BC = 7t$. According to the \href{https://artofproblemsolving.com/wiki/index.php?title=Law\_of\_cosines}{law of cosines},

 
                    \begin{figure}[H]
                        \centering
                        \includegraphics[width=0.5\textwidth]{Solution/AIME_I_1989_6_Solution_Figure_1.png}
                    \end{figure}
                    
 \begin{align*}(7t)^2 &= (8t)^2 + 100^2 - 2 \cdot 8t \cdot 100 \cdot \cos 60^\circ\\ 0 &= 15t^2 - 800t + 10000 = 3t^2 - 160t + 2000\\ t &= \frac{160 \pm \sqrt{160^2 - 4\cdot 3 \cdot 2000}}{6} = 20, \frac{100}{3}.\end{align*}


 Since we are looking for the earliest possible intersection, $20$ seconds are needed. Thus, $8 \cdot 20 = \boxed{160}$ meters is the solution.

 Alternatively, we can drop an altitude from $C$ and arrive at the same answer.

 





\textbf{Solution 2}


                    \begin{figure}[H]
                        \centering
                        \includegraphics[width=0.5\textwidth]{Solution/AIME_I_1989_6_Solution_Figure_2.png}
                    \end{figure}
                    Let $P$ be the point of intersection between the skaters, Allie and Billie. We can draw a line that goes through $P$ and is parallel to $\overline{AB}$. Letting this line be the $x$-axis, we can reflect $B$ over the $x$-axis to get $B'$. As reflections preserve length, $B'X = XB$.

 We then draw lines $BB'$ and $PB'$. We can let the foot of the perpendicular from $P$ to $BB'$ be $X$, and we can let the foot of the perpendicular from $P$ to $AB$ be $Y$. In doing so, we have constructed rectangle $PXBY$.

 By $d=rt$, we have $AP = 8t$ and $PB = PB' = 7t$, where $t$ is the number of seconds it takes the skaters to meet. Furthermore, we have $30-60-90$ triangle $PAY$, so $AY = 4t$, and $PY = 4t\sqrt{3}$. Since we have $PY = XB = B'X$, $B'X = 4t\sqrt{3}$. By Pythagoras, $PX = t$.

 As $PXBY$ is a rectangle, $PX = YB$. Thus $AY + YB = AB \Rightarrow AY + PX = AB$, so we get $4t + t = 100$. Solving for $t$, we find $t = 20$. 

 Our answer, $AP$, is equivalent to $8t$. Thus, $AP = 8 \cdot 20 = \boxed{160}$.

 





\textbf{Solution 3}

We can define $x$ to be the time elapsed since both Allie and Billie moved away from points $A$ and $B$ respectfully. Also, set the point of intersection to be $M$. Then we can produce the following diagram:

 
                    \begin{figure}[H]
                        \centering
                        \includegraphics[width=0.5\textwidth]{Solution/AIME_I_1989_6_Solution_Figure_3.png}
                    \end{figure}
                    

 Now, if we drop an altitude from point$M$, we get :

 
                    \begin{figure}[H]
                        \centering
                        \includegraphics[width=0.5\textwidth]{Solution/AIME_I_1989_6_Solution_Figure_4.png}
                    \end{figure}
                    

 We know this from the $30-60-90$ triangle that is formed. From this we get that: 
 \[(7x)^2 = (4 \sqrt{3} x)^2 + (100-4x)^2\]

 
 \[\Longrightarrow 49x^2 - 48x^2 = x^2 = (100-4x)^2\]

 
 \[\Longrightarrow 0=(100-4x)^2 - x^2 = (100-3x)(100-5x)\].

 Therefore, we get that $x = \frac{100}{3}$ or $x = 20$. Since $20< \frac{100}{3}$, we have that $x=20$ (since the problem asks for the quickest possible meeting point), so the distance Allie travels before meeting Billie would be $8 \cdot x = 8 \cdot 20 = \boxed{160}$ meters.

 





\textbf{Solution 4}


                    \begin{figure}[H]
                        \centering
                        \includegraphics[width=0.5\textwidth]{Solution/AIME_I_1989_6_Solution_Figure_5.png}
                    \end{figure}
                    

 Drop the altitude from $M$ to $AB$, and call it $P$. $\triangle AMP$ is a $30-60-90$ triangle, so $AP = 4t$ and $MP = 4\sqrt{3}t$, and by the Pythagorean theorem on $\triangle MPB$, $PB = t$. $AP + BP = AB$, so $t=20$. Therefore, $8t = \boxed{160}$.

 




	\item (\textit{AIME I 1989}) If the integer $k$ is added to each of the numbers $36$, $300$, and $596$, one obtains the squares of three consecutive terms of an arithmetic series. Find $k$.


 
 \textbf{Answer: }925



\textbf{Solution 1}

Call the terms of the \href{https://artofproblemsolving.com/wiki/index.php?title=Arithmetic\_progression}{arithmetic progression} $a,\ a + d,\ a + 2d$, making their squares $a^2,\ a^2 + 2ad + d^2,\ a^2 + 4ad + 4d^2$.

 We know that $a^2 = 36 + k$ and $(a + d)^2 = 300 + k$, and subtracting these two we get $264 = 2ad + d^2$ (1). Similarly, using $(a + d)^2 = 300 + k$ and $(a + 2d)^2 = 596 + k$, subtraction yields $296 = 2ad + 3d^2$ (2). 

 Subtracting the first equation from the second, we get $2d^2 = 32$, so $d = 4$. Substituting backwards yields that $a = 31$ and $k = \boxed{925}$.

 





\textbf{Solution 2 (Straighforward, but has big numbers)}

Since terms in an arithmetic progression have constant differences, 
 \[\sqrt{300+k}-\sqrt{36+k}=\sqrt{596+k}-\sqrt{300+k}\]
 \[\implies 2\sqrt{300+k} = \sqrt{596+k}+\sqrt{36+k}\]
 \[\implies 4k+1200=596+k+36+k+2\sqrt{(596+k)(36+k)}\]
 \[\implies 2k+568=2\sqrt{(596+k)(36+k)}\]
 \[\implies k+284=\sqrt{(596+k)(36+k)}\]
 \[\implies k^2+568k+80656=k^2+632k+21456\]
 \[\implies 568k+80656=632k+21456\]
 \[\implies 64k = 59200\]
 \[\implies k = \boxed{925}\]

 





\textbf{Solution 3}

Let the arithmetic sequence be $a-d$, $a$, and $a+d$. Then $(a+d)^2-a^2 = 296$, but using the difference of squares, $d(2a+d)=296$. Also, $a^2-(a-d)^2 = 264$, and using the difference of squares we get $d(2a-d) = 264$. Subtracting both equations gives $2d^2 = 32$, $d = 4$, and $a = 35$. Since $a = 35$, $a^2 = 1225 = 300+k$ and $k = \boxed{925}$.

 




	\item (\textit{AIME I 1989}) Assume that $x_1,x_2,\ldots,x_7$ are real numbers such that

 \begin{align*} x_1 + 4x_2 + 9x_3 + 16x_4 + 25x_5 + 36x_6 + 49x_7 &= 1, \\ 4x_1 + 9x_2 + 16x_3 + 25x_4 + 36x_5 + 49x_6 + 64x_7 &= 12, \\ 9x_1 + 16x_2 + 25x_3 + 36x_4 + 49x_5 + 64x_6 + 81x_7 &= 123. \end{align*}

Find the value of $16x_1+25x_2+36x_3+49x_4+64x_5+81x_6+100x_7$.


 
 \textbf{Answer: }334



\textbf{Solution 1 (Quadratic Function)}

Note that each given equation is of the form 
 \[f(k)=k^2x_1+(k+1)^2x_2+(k+2)^2x_3+(k+3)^2x_4+(k+4)^2x_5+(k+5)^2x_6+(k+6)^2x_7\] for some $k\in\{1,2,3\}.$

 When we expand $f(k)$ and combine like terms, we obtain a quadratic function of $k:$ 
 \[f(k)=ak^2+bk+c,\] where $a,b,$ and $c$ are linear combinations of $x_1,x_2,x_3,x_4,x_5,x_6,$ and $x_7.$ 

 We are given that
 \begin{alignat*}{10} f(1)&=\phantom{42}a+b+c&&=1, \\ f(2)&=4a+2b+c&&=12, \\ f(3)&=9a+3b+c&&=123, \end{alignat*}and we wish to find $f(4).$

 We eliminate $c$ by subtracting the first equation from the second, then subtracting the second equation from the third:
 \begin{align*} 3a+b&=11, \\ 5a+b&=111. \end{align*}
By either substitution or elimination, we get $a=50$ and $b=-139.$ Substituting these back produces $c=90.$

 Finally, the answer is 
 \[f(4)=16a+4b+c=\boxed{334}.\]

 





\textbf{Solution 2 (Linear Combination)}

For simplicity purposes, we number the given equations $(1),(2),$ and $(3),$ in that order. Let 
 \[16x_1+25x_2+36x_3+49x_4+64x_5+81x_6+100x_7=S. \hspace{29.5mm}(4)\]Subtracting $(1)$ from $(2),$ subtracting $(2)$ from $(3),$ and subtracting $(3)$ from $(4),$ we obtain the following equations, respectively:
 \begin{align*} 3x_1 + 5x_2 +  7x_3 +  9x_4 + 11x_5 + 13x_6 + 15x_7 &=11, \hspace{20mm}&(5) \\ 5x_1 + 7x_2 +  9x_3 + 11x_4 + 13x_5 + 15x_6 + 17x_7 &=111, &(6) \\ 7x_1 + 9x_2 + 11x_3 + 13x_4 + 15x_5 + 17x_6 + 19x_7 &=S-123. &(7) \\ \end{align*}
Subtracting $(5)$ from $(6)$ and subtracting $(6)$ from $(7),$ we obtain the following equations, respectively:
 \begin{align*} 2x_1+2x_2+2x_3+2x_4+2x_5+2x_6+2x_7&=100, &(8) \\ 2x_1+2x_2+2x_3+2x_4+2x_5+2x_6+2x_7&=S-234. \hspace{20mm}&(9) \end{align*}
Finally, applying the Transitive Property to $(8)$ and $(9)$ gives $S-234=100,$ from which $S=\boxed{334}.$

 





\textbf{Solution 3 (Finite Differences by Arithmetic)}

Note that the second differences of all quadratic sequences must be constant (but nonzero). One example is the following sequence of perfect squares:

 
                    \begin{figure}[H]
                        \centering
                        \includegraphics[width=0.5\textwidth]{Solution/AIME_I_1989_8_Solution_Figure_1.png}
                    \end{figure}
                    

 Label equations $(1),(2),(3),$ and $(4)$ as Solution 2 does. Since the coefficients of $x_1,x_2,x_3,x_4,x_5,x_6,x_7,$ or $(1,4,9,16),(4,9,16,25),(9,16,25,36),(16,25,36,49),(25,36,49,64),(36,49,64,81),(49,64,81,100),$ respectively, all form quadratic sequences with second differences $2,$ we conclude that the second differences of equations $(1),(2),(3),(4)$ must be constant.

 It follows that the second differences of $(1,12,123,S)$ must be constant, as shown below:

 
                    \begin{figure}[H]
                        \centering
                        \includegraphics[width=0.5\textwidth]{Solution/AIME_I_1989_8_Solution_Figure_2.png}
                    \end{figure}
                    

 Finally, we have $d_2=100,$ from which
 \begin{align*} S&=123+d_1 \\ &=123+(111+d_2) \\ &=\boxed{334}. \end{align*}
~MRENTHUSIASM

 





\textbf{Solution 4 (Finite Differences by Algebra)}

Notice that we may rewrite the equations in the more compact form as:
 \begin{align*} \sum_{i=1}^{7}i^2x_i&=c_1, \\ \sum_{i=1}^{7}(i+1)^2x_i&=c_2, \\ \sum_{i=1}^{7}(i+2)^2x_i&=c_3, \\ \sum_{i=1}^{7}(i+3)^2x_i&=c_4, \end{align*}
where $c_1=1, c_2=12, c_3=123,$ and $c_4$ is what we are trying to find.

 Now consider the polynomial given by $f(z) = \sum_{i=1}^7 (i+z)^2x_i$ (we are only treating the $x_i$ as coefficients).

 Notice that $f$ is in fact a quadratic. We are given $f(0)=c_1, f(1)=c_2, f(2)=c_3$ and are asked to find $f(3)=c_4$. Using the concept of finite differences (a prototype of differentiation) we find that the second differences of consecutive values is constant, so that by arithmetic operations we find $c_4=\boxed{334}$. 

 Alternatively, applying finite differences, one obtains 
 \[c_4 = {3 \choose 2}f(2) - {3 \choose 1}f(1) + {3 \choose 0}f(0) =334.\]

 





\textbf{Solution 5 (Assumption)}

The idea is to multiply the first, second and third equations by $a,b,$ and $c,$ respectively.

 We can only consider the coefficients of $x_1,x_2,$ and $x_3:$
 \begin{align} a+4b+9c&=16, \\ 4a+9b+16c&=25, \\ 9a+16b+25c&=36. \end{align}Subtracting $(1)$ from $(2),$ we get 
 \[3a+5b+7c=9. \hspace{15mm}(4)\]Subtracting $3\cdot(4)$ from $(3),$ we get 
 \[b+4c=9. \hspace{25.5mm}(5)\]Subtracting $(1)$ from $4\cdot(5),$ we get 
 \[7c-a=20. \hspace{23mm}(6)\]From $(5)$ and $(6),$ we have $(a,b,c)=(7c-20,9-4c,c).$ Substituting this into $(2)$ gives $(a,b,c)=(1,-3,3).$

 Therefore, the answer is $1\cdot1+12\cdot(-3) + 123\cdot3 = \boxed{334}.$

 





\textbf{Solution 6 (Assumption)}

We let $(x_4,x_5,x_6,x_7)=(0,0,0,0)$. Thus, we have 
 \begin{align*}  x_1+4x_2+9x_3&=1,\\ 4x_1+9x_2+16x_3&=12,\\ 9x_1+16x_2+25x_3&=123.\\  \end{align*}
Grinding this out, we have $(x_1,x_2,x_3)=\left(\frac{797}{4},-229,\frac{319}{4}\right)$ which gives $\boxed{334}$ as our final answer. 

 





\textbf{Solution 7 (Similar to Solutions 3 and 4)}

Let $s_n = n^2$ be the sequence of perfect squares.By either expanding or via finite differences, one can prove the miraculous recursion
 \[s_n = 3s_{n-1} - 3s_{n-2} + s_{n-3}.\]Hence, the answer is simply
 \[3 \cdot 123 - 3 \cdot 12 + 1 = \boxed{334}.\]

 Note: You can prove this recursion through considering the coefficients for each $x_i$ as $a^2, (a+1)^2, (a+2)^2$ for the equations respectively.

 All we are left to do is find $(a+3)^2$ in terms of the above $3$ quadratics. Because we have an $a^2$ as one, we don't have to worry about the $a^2$ term, we can just manipulate it by adding or subtracting any number of the first quadratic.

 We are left to find a combination of the second and third quadratics that have the linear term and constant of $(a+3)^2$. With little experimentation we find that $3(a+2)^2 - 3(a+1)^2 = 6a+9$. All we have to add back is one $a^2$.

 Thus, the recursion holds. That for any term $x_i$ the coefficients would be recursible by $a^2$ = $3(a-1)^2 - 3(a-2)^2 + (a-3)$(Where we have set $a$ as $a-3$)

 





\textbf{Solution 8}

We number our equations with $(n)$ denoting equation $n$
 \begin{align*} x_1 + 4x_2 + 9x_3 + 16x_4 + 25x_5 + 36x_6 + 49x_7 &= 1  &(1)\\ 4x_1 + 9x_2 + 16x_3 + 25x_4 + 36x_5 + 49x_6 + 64x_7 &= 12  &(2) \\ 9x_1 + 16x_2 + 25x_3 + 36x_4 + 49x_5 + 64x_6 + 81x_7 &= 123   &(3) \end{align*}
See that we want to find the value of 
 \[7x_{1}+9x_{2}+11x_{3}+13x_{4}+15x_{5}+17x_{6}+19x_{7}\] and add that to $(3)$. Keep this in mind. 

 
Subtracting $(1)$ from $(2)$ yields
 \[3x_{1}+5x_{2}+7x_{3}+9x_{4}+11x_{5}+13x_{6}+15x_{7} = 11 \ \ (4)\]and subtracting $(2)$ from $(3)$ yields
 \[5x_{1}+7x_{2}+9x_{3}+11x_{4}+13x_{5}+15x_{6}+17x_{7} = 111 \ \ \ (5)\]subtracting $(4)$ from $(5)$ yields
 \[2x_{1}+2x_{2}+2x_{3}+2x_{4}+2x_{5}+2x_{6}+2x_{7} = 100 \ \ \ (6)\]Adding this back to $(5)$ yields
 \[7x_{1}+9x_{2}+11x_{3}+13x_{4}+15x_{5}+17x_{6}+19x_{7}=211 \ \ \ (7)\]Now, simply adding $(7)$ to $(3)$ yields
 \[16x_1+25x_2+36x_3+49x_4+64x_5+81x_6+100x_7 = \boxed{334}\]~shockfront99(omg this is my first AIME solve)

 




	\item (\textit{AIME I 1989}) One of Euler's conjectures was disproved in the 1960s by three American mathematicians when they showed there was a positive integer such that 
 \[133^5+110^5+84^5+27^5=n^{5}.\] Find the value of $n$.


 
 \textbf{Answer: }144



\textbf{Solution 1 (FLT, CRT, Inequalities)}

Taking the given equation modulo $2,3,$ and $5,$ respectively, we have
 \begin{align*} n^5&\equiv0\pmod{2}, \\ n^5&\equiv0\pmod{3}, \\ n^5&\equiv4\pmod{5}. \end{align*}
By either \textbf{Fermat's Little Theorem (FLT)} or inspection, we get
 \begin{align*} n&\equiv0\pmod{2}, \\ n&\equiv0\pmod{3}, \\ n&\equiv4\pmod{5}. \end{align*}
By either the \textbf{Chinese Remainder Theorem (CRT)} or inspection, we get $n\equiv24\pmod{30}.$

 It is clear that $n>133,$ so the possible values for $n$ are $144,174,204,\ldots.$ Note that
 \begin{align*} n^5&=133^5+110^5+84^5+27^5 \\ &<133^5+110^5+(84+27)^5 \\ &=133^5+110^5+111^5 \\ &<3\cdot133^5, \end{align*}
from which $\left(\frac{n}{133}\right)^5<3.$

 If $n\geq174,$ then
 \begin{align*} \left(\frac{n}{133}\right)^5&>1.3^5 \\ &=1.3^2\cdot1.3^2\cdot1.3 \\ &>1.6\cdot1.6\cdot1.3 \\ &=2.56\cdot1.3 \\ &>2.5\cdot1.2 \\ &=3, \end{align*}
which arrives at a contradiction. Therefore, we conclude that $n=\boxed{144}.$

 





\textbf{Solution 2}

Note that $n$ is even, since the LHS consists of two odd and two even numbers. By Fermat's Little Theorem, we know $n^5\equiv n\pmod{5}.$ Hence, 
 \[n\equiv3+0+4+2\equiv4\pmod{5}.\]Continuing, we examine the equation modulo $3,$ 
 \[n\equiv1-1+0+0\equiv0\pmod{3}.\]Thus, $n$ is divisible by three and leaves a remainder of four when divided by $5.$ It is obvious that $n>133,$ so the only possibilities are $n = 144$ or $n \geq 174.$ It quickly becomes apparent that $174$ is much too large, so $n$ must be $\boxed{144}.$

 





\textbf{Solution 3}

We can cheat a little bit and approximate, since we are dealing with such large numbers. As above, $n^5\equiv n\pmod{5},$ and it is easy to see that $n^5\equiv n\pmod 2.$ Therefore, $133^5+110^5+84^5+27^5\equiv 3+0+4+7\equiv 4\pmod{10},$ so the last digit of $n$ is $4.$

 We notice that $133,110,84,$ and $27$ are all very close or equal to multiples of $27.$ We can rewrite $n^5$ as approximately equal to $27^5(5^5+4^5+3^5+1^5) = 27^5(4393).$ This means $\frac{n^5}{27^5}$ must be close to $4393.$

 Note that $134$ will obviously be too small, so we try $144$ and get $\left(\frac{144}{27}\right)^5=\left(\frac{16}{3}\right)^5.$ Bashing through the division, we find that $\frac{1048576}{243}\approx 4315,$ which is very close to $4393.$ It is clear that $154$ will not give any closer of an answer, given the rate that fifth powers grow, so we can safely assume that $\boxed{144}$ is the answer.

 





\textbf{Solution 4}

In this solution we take advantage of the large numbers and utilize parity properties to give us a very good guess at the answer. The units digits of $133^5, 110^5, 84^5, 27^5$ are $3, 0, 4, 7,$ respectively, so the units digit of $n^5$ is $4.$ This tells us $n$ is even. Since we are dealing with enormous numbers, $n$ should not be that far from $133.$ Note that $n$'s units digit is $0, 2, 4, 6,$ or $8.$ When to the power of $5,$ they each give $0, 2, 4, 6,$ and $8$ as the units digits. This further clues us that $n$ ends in $4.$ 

 Clearly, $n>133,$ so we start with $134.$ Now we need a way of distinguishing between numbers with units digit $4.$ We can do this by finding the last three digits for each of $133^5, 110^5, 84^5,$ and $27^5,$ which is not that difficult. For $133^5,$ we have
 \begin{align*} 133^5&=133^2\cdot133^2\cdot133 \\ &\equiv689\cdot689\cdot133 \\ &\equiv721\cdot133 \\ &\equiv893\pmod{1000}. \end{align*}
By the same reasoning, we get
 \begin{align*} n^5&=133^5+110^5+84^5+27^5 \\ &\equiv893+0+424+907 \\ &\equiv224\pmod{1000}. \end{align*}
Note that
 \begin{align*} 134^5&\equiv424\pmod{1000}, \\ 144^5&\equiv224\pmod{1000}, \\ 154^5&\equiv024\pmod{1000}, \\ 164^5&\equiv824\pmod{1000}, \\ 174^5&\equiv624\pmod{1000}, \\ 184^5&\equiv424\pmod{1000}, \\ 194^5&\equiv224\pmod{1000}. \end{align*}
By observations, $n=194$ is obviously an overestimate. So, the answer is $n=\boxed{144}.$

 





\textbf{Solution 5}

First, we take mod $2$ on both sides to get $n^5\equiv 0\pmod{2}\implies n\equiv 0\pmod{2}$. Mod $3$ gives $n^5\equiv 0\pmod{3}\implies n\equiv 0\pmod{3}$. Also, mod $5$ gives $n^5\equiv -1\pmod{5}\implies n\equiv -1\pmod{5}$ (by FLT). Finally, note that mod $7$ gives $n^5\equiv 2\pmod{7}\implies n^{-1}\equiv 2\pmod{7}\implies n\equiv 4\pmod{7}$. Thus,
 \begin{align*} n&\equiv 0\pmod{2}, \\ n&\equiv 0\pmod{3}, \\ n&\equiv -1\pmod{5}, \\ n&\equiv 4\pmod{7}. \end{align*}
By CRT, $n\equiv 144\pmod{210}$, so $n$ is one of $144, 354, ...$. However, $133^5 + 110^5 + 84^5 + 27^5 < 4\cdot 133^5 < (2\cdot 133)^5 < 354^5$, so $n < 354$. Thus, $n = \boxed{144}$.

 





\textbf{Solution 6 (Brute Force)}

We have 
 \[n^5 = 133^5 + 110^5 + 84^5 +27^5 = 61917364224,\] for which $n = \sqrt [5]{61917364224} = \boxed{144}.$

 





\textbf{Solution 7 (Cheese)}

By approximation, $110 \approx \frac{5}{6}\cdot133$, $84 \approx \frac{7}{11}\cdot133$, and $23 \approx \frac{1}{5}\cdot133$. Adding the fifth power of each fraction of $133$ gives a total of roughly $1.506\cdot133^5$. By testing, we know that the fifth root of this value is less than $1.1\cdot133=146.3$, as $1.1^5=1.611$, meaning we have reduced $n$ to $133\le{n}\le{147}$.The last digit of all four numbers, in order, are $3$, $0$, $4$, and $7$, which is gained through writing out the cycles of the last digit up to the fifth power. This means that the answer has a units digit of $4$, reducing $n$ down to either $134$ or $144$. Since the previous cheesing shows obviously that the value is closer to $146.3$ than it is to $133$, we know our final answer of $\boxed{144}$ is correct.

 Cheese solution by \href{https://artofproblemsolving.com/wiki/index.php?title=User:Juwushu}{juwushu}.

 





\textbf{Solution 8 (Guesswork or Bash)}

We know that $a^5 \equiv a \pmod{10}$ for all a. You can prove this later. So we know the number ends with 4. Since 133 is the biggest number, while everything else is small, we can safely assume that it is $\boxed{144}$, because 5th powers vary wildly. It is not 134 because $110^5$ plays a role in the sum.  We can also just take $130^5$ and the other numbers $110^5, 80^5,20^5$, add them up, and we realize that $140^5$ is close to this. Then we also get $\boxed{144}$ 

 





\textbf{Solution 9 (Guess, only if you have no time)}

The number should be around 140, and since most solutions in the AIME are 'round' numbers, we can guess 144.~Sigma\_Bob

 \textbf{Remark} We could also take this expression mod $7,$ which would give $n\equiv4\pmod{7}.$ Then the only reasonable number that works is $144,$ so we could safely assume this is the answer.

 




	\item (\textit{AIME I 1989}) Let $a$, $b$, $c$ be the three sides of a triangle, and let $\alpha$, $\beta$, $\gamma$, be the angles opposite them. If $a^2+b^2=1989c^2$, find


 $\frac{\cot \gamma}{\cot \alpha+\cot \beta}$
 \textbf{Answer: }994



\textbf{Solution 1}

We draw the \href{https://artofproblemsolving.com/wiki/index.php?title=Altitude}{altitude} $h$ to $c$, to get two \href{https://artofproblemsolving.com/wiki/index.php?title=Right\_triangle}{right triangles}.

 
                    \begin{figure}[H]
                        \centering
                        \includegraphics[width=0.5\textwidth]{Solution/AIME_I_1989_10_Solution_Figure_1.png}
                    \end{figure}
                    Then $\cot{\alpha}+\cot{\beta}=\frac{c}{h}$, from the definition of the \href{https://artofproblemsolving.com/wiki/index.php?title=Cotangent}{cotangent}.

 Let $K$ be the area of $\triangle ABC.$ Then $h=\frac{2K}{c}$, so $\cot{\alpha}+\cot{\beta}=\frac{c^2}{2K}$.

 By identical logic, we can find similar expressions for the sums of the other two cotangents:
 \begin{align*} \cot \alpha + \cot \beta &= \frac{c^2}{2K} \\ \cot \beta + \cot \gamma &= \frac{a^2}{2K} \\ \cot \gamma + \cot \alpha &= \frac{b^2}{2K}. \end{align*}
Adding the last two equations, subtracting the first, and dividing by 2, we get
 \[\cot \gamma = \frac{a^2 + b^2 - c^2}{4K}.\]Therefore
 \begin{align*} \frac{\cot \gamma}{\cot \alpha + \cot \beta} &= \frac{(a^2 + b^2 - c^2)/(4K)}{c^2/(2K)} \\ &= \frac{a^2 + b^2 - c^2}{2c^2} \\ &= \frac{1989 c^2 - c^2}{2c^2} \\ &= \frac{1988}{2} = \boxed{994}. \end{align*}


 





\textbf{Solution 2}

By the \href{https://artofproblemsolving.com/wiki/index.php?title=Law\_of\_cosines}{law of cosines},
 \[\cos \gamma = \frac{a^2 + b^2 - c^2}{2ab}.\] So, by the extended \href{https://artofproblemsolving.com/wiki/index.php?title=Law\_of\_sines}{law of sines},
 \[\cot \gamma = \frac{\cos \gamma}{\sin \gamma} = \frac{a^2 + b^2 - c^2}{2ab} \cdot \frac{2R}{c} = \frac{R}{abc} (a^2 + b^2 - c^2).\]Identical logic works for the other two angles in the triangle.  So, the cotangent of any angle in the triangle is directly proportional to the sum of the squares of the two adjacent sides, minus the square of the opposite side.  Therefore
 \[\frac{\cot \gamma}{\cot \alpha + \cot \beta} = \frac{a^2 + b^2 - c^2}{(-a^2 + b^2 + c^2) + (a^2 - b^2 + c^2)} = \frac{a^2 + b^2 - c^2}{2c^2}.\]We can then finish as in solution 1.

 





\textbf{Solution 3}

We start as in solution 1, though we'll write $A$ instead of $K$ for the area.  Now we evaluate the numerator:

 
 \[\cot{\gamma}=\frac{\cos{\gamma}}{\sin{\gamma}}\]

 From the \href{https://artofproblemsolving.com/wiki/index.php?title=Law\_of\_Cosines}{Law of Cosines} and the sine area formula, 

 
 \begin{align*}\cos{\gamma}&=\frac{1988c^2}{2ab}\\ \sin{\gamma}&= \frac{2A}{ab}\\ \cot{\gamma}&= \frac{\cos \gamma}{\sin \gamma} = \frac{1988c^2}{4A} \end{align*}


 Then $\frac{\cot \gamma}{\cot \alpha+\cot \beta}=\frac{\frac{1988c^2}{4A}}{\frac{c^2}{2A}}=\frac{1988}{2}=\boxed{994}$.

 





\textbf{Solution 4}


 \begin{align*} \cot{\alpha} + \cot{\beta} &= \frac {\cos{\alpha}}{\sin{\alpha}} + \frac {\cos{\beta}}{\sin{\beta}} = \frac {\sin{\alpha}\cos{\beta} + \cos{\alpha}\sin{\beta}}{\sin{\alpha}\sin{\beta}}\\ &= \frac {\sin{(\alpha + \beta)}}{\sin{\alpha}\sin{\beta}} = \frac {\sin{\gamma}}{\sin{\alpha}\sin{\beta}} \end{align*}


 By the Law of Cosines,

 
 \[a^2 + b^2 - 2ab\cos{\gamma} = c^2 = 1989c^2 - 2ab\cos{\gamma} \implies ab\cos{\gamma} = 994c^2\]

 Now

 
 \begin{align*}\frac {\cot{\gamma}}{\cot{\alpha} + \cot{\beta}} &= \frac {\cot{\gamma}\sin{\alpha}\sin{\beta}}{\sin{\gamma}} = \frac {\cos{\gamma}\sin{\alpha}\sin{\beta}}{\sin^2{\gamma}} = \frac {ab}{c^2}\cos{\gamma} = \frac {ab}{c^2} \cdot \frac {994c^2}{ab}\\ &= \boxed{994}\end{align*}


 


 





\textbf{Solution 5}

Use Law of cosines to give us $c^2=a^2+b^2-2ab\cos(\gamma)$ or therefore $\cos(\gamma)=\frac{994c^2}{ab}$.  Next, we are going to put all the sin's in term of $\sin(a)$.  We get $\sin(\gamma)=\frac{c\sin(a)}{a}$.  Therefore, we get $\cot(\gamma)=\frac{994c}{b\sin a}$.

 Next, use Law of Cosines to give us $b^2=a^2+c^2-2ac\cos(\beta)$.  Therefore, $\cos(\beta)=\frac{a^2-994c^2}{ac}$.  Also, $\sin(\beta)=\frac{b\sin(a)}{a}$.  Hence, $\cot(\beta)=\frac{a^2-994c^2}{bc\sin(a)}$.  

 Lastly, $\cos(\alpha)=\frac{b^2-994c^2}{bc}$. Therefore, we get $\cot(\alpha)=\frac{b^2-994c^2}{bc\sin(a)}$.

 Now, $\frac{\cot(\gamma)}{\cot(\beta)+\cot(\alpha)}=\frac{\frac{994c}{b\sin a}}{\frac{a^2-994c^2+b^2-994c^2}{bc\sin(a)}}$.  After using $a^2+b^2=1989c^2$, we get $\frac{994c*bc\sin a}{c^2b\sin a}=\boxed{994}$.

 


 





\textbf{Solution 6}

Let $\gamma$ be $(180-\alpha-\beta)$

 $\frac{\cot \gamma}{\cot \alpha+\cot \beta} = \frac{\frac{-\tan \alpha \tan \beta}{\tan(\alpha+\beta)}}{\tan \alpha + \tan \beta} = \frac{(\tan \alpha \tan \beta)^2-\tan \alpha \tan \beta}{\tan^2 \alpha + 2\tan \alpha \tan \beta +\tan^2 \beta}$

 WLOG, assume that $a$ and $c$ are legs of right triangle $abc$ with $\beta = 90^o$ and $c=1$

 By the Pythagorean theorem, we have $b^2=a^2+1$, and the given $a^2+b^2=1989$. Solving the equations gives us $a=\sqrt{994}$ and $b=\sqrt{995}$. We see that $\tan \beta = \infty$, and $\tan \alpha = \sqrt{994}$. 

 Our derived equation equals $\tan^2 \alpha$ as $\tan \beta$ approaches infinity.Evaluating $\tan^2 \alpha$, we get $\boxed{994}$.

 


 





\textbf{Solution 7}

As in Solution 1, drop an altitude $h$ to $c$. Let $h$ meet $c$ at $P$, and let $AP = x, BP = y$.

 
                    \begin{figure}[H]
                        \centering
                        \includegraphics[width=0.5\textwidth]{Solution/AIME_I_1989_10_Solution_Figure_2.png}
                    \end{figure}
                    Then, $\cot{\alpha} = \frac{1}{\tan{\alpha}} = \frac{x}{h}$, $\cot{\beta} = \frac{1}{\tan{\beta}} = \frac{y}{h}$. We can calculate $\cot{\gamma}$ using the \href{https://artofproblemsolving.com/wiki/index.php?title=Tangent\_addition\_formula}{tangent addition formula}, after noticing that $\cot{\gamma} = \frac{1}{\tan{\gamma}}$. So, we find that \begin{align*}\cot{\gamma} &= \frac{1}{\tan{\gamma}} \\&= \frac{1}{\frac{\frac{x}{h} + \frac{y}{h}}{1 - \frac{xy}{h^2}}} \\&= \frac{1}{\frac{(x+y)h}{h^2 - xy}} \\&= \frac{h^2 - xy}{(x+y)h}.\end{align*}


 So now we can simplify our original expression:\begin{align*}\frac{\cot{\gamma}}{\cot{\alpha} + \cot{\beta}} &= \frac{\frac{h^2 - xy}{(x+y)h}}{\frac{x + y}{h}} \\&= \frac{h^2 - xy}{(x+y)^2}.\end{align*}


 But notice that $x+y = c$, so this becomes 
 \[\frac{h^2 - xy}{c^2}.\]Now note that we can use the \href{https://artofproblemsolving.com/wiki/index.php?title=Pythagorean\_theorem}{Pythagorean theorem} to calculate $h^2$, we get that 
 \[h^2 = \frac{a^2 - y^2 + b^2 - x^2}{2}.\]So our expression simplifies to 
 \[\frac{1989c^2 - (x+y)^2}{2c^2}\]since $a^2 + b^2 = 1989c^2$ from the problem and that there is another $-\frac{2xy}{2}$ after the $h^2$ in our expression. Again note that $x+y = c$, so it again simplifies to $\frac{1988c^2}{2c^2}$, or $\boxed{994}$.

 


 





\textbf{Solution 8 (Quick and Easy)}

Since no additional information is given, we can assume that triangle ABC is right with the right angle at B.We can use the \href{https://artofproblemsolving.com/wiki/index.php?title=Pythagorean\_theorem}{Pythagorean theorem} to say 
 \[c^2+a^2=b^2\] We can now solve for $a$ in terms of $c$

 
 \[c^2+a^2=1989c^2-a^2\]
 \[a^2=994c^2\]
 \[a=\sqrt{994}c\]

 Using the definition of cotangent

 
 \[cot(A)=\frac{c}{a}=\frac{1}{\sqrt{994}}\]
 \[cot(B)=cot(90)=0\]
 \[cot(C)=\frac{a}{c}=\sqrt{994}\]Plugging into our desired expression, we get $\boxed{994}$

 





\textbf{Solution 9 (Quick and Easy), similar idea to Solution 8}

Assume the triangle is isosceles and $c=1$.  Since $a=b$, $2a^2=1989$, $\rightarrow a=b=\sqrt{\frac{1989}{2}}$

 Since this is an isosceles triangle, we can find $\cot(\alpha)$, $\cot(\beta)$, by just drawing the height.  Via Pythagorean Theorem, the height is $\frac{\sqrt{3977}}{2}$. 

 Therefore we find that $\tan(\alpha) = \tan(\beta) = \sqrt{3977}$.  $\tan(\gamma) = \tan(180-(\alpha+\beta)) = -\tan(\alpha+\beta)$.  Via tangent sum formula, $\tan(\alpha+\beta) = -\frac{\sqrt{3977}}{1988}$, so $\tan(\gamma) = \frac{\sqrt{3977}}{1988}$.

 Since $\cot(\theta) = \frac{1}{\tan(\theta)}$,  $\cot(\alpha) = \frac{1}{\sqrt{3977}}$, $\cot(\beta) = \frac{1}{\sqrt{3977}}$, $\cot(\gamma) = \frac{1988}{\sqrt{3977}}$.

 Plugging this in, almost everything cancels and we get $\frac{1988}{2} = \boxed{994}$.

 




	\item (\textit{AIME I 1989}) A sample of 121 integers is given, each between 1 and 1000 inclusive, with repetitions allowed. The sample has a unique mode (most frequent value). Let $D^{}_{}$ be the difference between the mode and the arithmetic mean of the sample. What is the largest possible value of $\lfloor D^{}_{}\rfloor$? (For real $x^{}_{}$, $\lfloor x^{}_{}\rfloor$ is the greatest integer less than or equal to $x^{}_{}$.)


 
 \textbf{Answer: }947



\textbf{Solution 1}

Let the mode be $x$, which we let appear $n > 1$ times. We let the arithmetic mean be $M$, and the sum of the numbers $\neq x$ be $S$. Then 
 \begin{align*} D &= \left|M-x\right| = \left|\frac{S+xn}{121}-x\right| = \left|\frac{S}{121}-\left(\frac{121-n}{121}\right)x\right| \end{align*}
As $S$ is essentially independent of $x$, it follows that we wish to minimize or maximize $x$ (in other words, $x \in \{1,1000\}$). Indeed, $D(x)$ is symmetric about $x = 500.5$; consider replacing all of numbers $x_i$ in the sample with $1001-x_i$, and the value of $D$ remains the same. So, \href{https://artofproblemsolving.com/wiki/index.php?title=Without\_loss\_of\_generality}{without loss of generality}, let $x=1$. Now, we would like to maximize the quantity 

 $\frac{S}{121}-\left(\frac{121-n}{121}\right)(1) = \frac{S+n}{121}-1$$S$ contains $121-n$ numbers that may appear at most $n-1$ times. Therefore, to maximize $S$, we would have $1000$ appear $n-1$ times, $999$ appear $n-1$ times, and so forth. We can thereby represent $S$ as the sum of $n-1$ arithmetic series of $1000, 999, \ldots, 1001 - \left\lfloor \frac{121-n}{n-1} \right\rfloor$. We let $k = \left\lfloor \frac{121-n}{n-1} \right\rfloor$, so

 $S = (n-1)\left[\frac{k(1000 + 1001 - k)}{2}\right] + R(n)$where $R(n)$ denotes the sum of the remaining $121-(n-1)k$ numbers, namely $R(n) = (121-(n-1)k)(1000-k)$. 

 At this point, we introduce the crude estimate<sup>[\href{https://artofproblemsolving.com\#note\_1}{1}]</sup> that $k=\frac{121-n}{n-1}$, so $R(n) = 0$ and
 \begin{align*}2S+2n &= (121-n)\left(2001-\frac{121-n}{n-1}\right)+2n = (120-(n-1))\left(2002-\frac{120}{n-1}\right)\end{align*}
Expanding (ignoring the constants, as these do not affect which $n$ yields a maximum) and scaling, we wish to minimize the expression $5(n-1) + \frac{36}{n-1}$. By \href{https://artofproblemsolving.com/wiki/index.php?title=AM-GM}{AM-GM}, we have $5(n-1) + \frac{36}{n-1} \ge 2\sqrt{5(n-1) \cdot \frac{36}{n-1}}$, with equality coming when $5(n-1) = \frac{36}{n-1}$, so $n-1 \approx 3$. Substituting this result and some arithmetic gives an answer of $\boxed{947}$.

 In less formal language, it quickly becomes clear after some trial and error that in our sample, there will be $n$ values equal to one and $n-1$ values each of $1000, 999, 998 \ldots$. It is fairly easy to find the \href{https://artofproblemsolving.com/wiki/index.php?title=Maximum}{maximum}. Try $n=2$, which yields $924$, $n=3$, which yields $942$,  $n=4$, which yields $947$, and $n=5$, which yields $944$. The maximum difference occurred at $n=4$, so the answer is $947$.

 





\textbf{Solution 2}

With the same reasoning as Solution 1, in order to get largest possible value of D, we can construct that our set of numbers as $\underbrace{1,1,1...1,}_\text{n times}\underbrace{2,2,2...2,}_\text{n times}\underbrace{3,3,3...3,}_\text{n times}........\underbrace{1000,1000,1000....}_\text{n+1  times}$ And, we need to find the value of n that makes the sum as low as possible. And, we can create a formula to make it easier. It isn't hard to find the sum. The numbers which are not 1000, average to $\frac{120}{2n}$ or $\frac{60}{n}$, and there are $120-n$ of them. So, they sum to $\frac{60}{n}(120-n)$. And, the sum of the numbers that are 1000 is $1000(n+1)$ so, our total sum gets us $1000(n+1)+120/2n(120-n)$ We want to minimize it, since the mode will always be 1000. And, testing the values n = 1, n = 2, n = 3, n = 4, we get these results.

 $n = 1: 2000+60*119 = 9140$

 $n = 2: 3000+30*118 = 6540$

 $n = 3: 4000+20*117 = 6340$

 $n = 4: 5000+15*116 = 6740$

 And, as n grows larger and larger from 4, the values will start increasing. Thus, the lowest possible sum is 6340. Dividing by 121, the lowest possible mean is 52.396...., and thus, the highest possible value of $D$ is 947.604, and the floor of that is $\boxed{947}$

 - AlexLikeMath

 \textbf{Notes} 
                \begin{itemize}
                    \item \href{https://artofproblemsolving.com\#ref\_1}{\textasciicircum{}} In fact, when $n = 2,3,4,5$ (which some simple testing shows that the maximum will occur around), it turns out that $\frac{121-n}{n-1}$ is an integer anyway, so indeed $k = \left\lfloor \frac{121-n}{n-1} \right\rfloor = \frac{121-n}{n-1}$.
                \end{itemize}


                




	\item (\textit{AIME I 1989}) Let $ABCD$ be a tetrahedron with $AB=41$, $AC=7$, $AD=18$, $BC=36$, $BD=27$, and $CD=13$, as shown in the figure. Let $d$ be the distance between the midpoints of edges $AB$ and $CD$. Find $d^{2}$.


 
                    \begin{figure}[H]
                        \centering
                        \includegraphics[width=0.5\textwidth]{Problem/AIME_I_1989_12_Problem_Figure_1.png}
                    \end{figure}
                    
 \textbf{Answer: }137



\textbf{Solution}

Call the midpoint of $\overline{AB}$ $M$ and the midpoint of $\overline{CD}$ $N$. $d$ is the \href{https://artofproblemsolving.com/wiki/index.php?title=Median}{median} of triangle $\triangle CDM$. The formula for the length of a median is $m=\sqrt{\frac{2a^2+2b^2-c^2}{4}}$, where $a$, $b$, and $c$ are the side lengths of triangle, and $c$ is the side that is bisected by median $m$. The formula is a direct result of the \href{https://artofproblemsolving.com/wiki/index.php?title=Law\_of\_Cosines}{Law of Cosines} applied twice with the angles formed by the median (\href{https://artofproblemsolving.com/wiki/index.php?title=Stewart\%27s\_Theorem}{Stewart's Theorem}). We can also get this formula from the parallelogram law, that the sum of the squares of the diagonals is equal to the squares of the sides of a parallelogram (\href{https://en.wikipedia.org/wiki/Parallelogram\_law}{https://en.wikipedia.org/wiki/Parallelogram\_law}).

 We first find $CM$, which is the median of $\triangle CAB$.

 
 \[CM=\sqrt{\frac{98+2592-1681}{4}}=\frac{\sqrt{1009}}{2}\]

 Now we must find $DM$, which is the median of $\triangle DAB$.

 
 \[DM=\frac{\sqrt{425}}{2}\]

 Now that we know the sides of $\triangle CDM$, we proceed to find the length of $d$.

 
 \[d=\frac{\sqrt{548}}{2} \Longrightarrow d^2=\frac{548}{4}=\boxed{137}\]

 





\textbf{Solution 2}

Let M be the midpoint of AB.DM is a median of triangle ABD and CM is a median of triangle ABCWe find that CM\textasciicircum{}2 = 1009/4 and DM\textasciicircum{}2 = 425/4We can then find the median of triangle CDM squared from the 3 side lengths of the triangle which gives us 137.Therefore \(d^2 = \boxed{137}.\)

 




	\item (\textit{AIME I 1989}) Let $S^{}_{}$ be a subset of $\{1,2,3^{}_{},\ldots,1989\}$ such that no two members of $S^{}_{}$ differ by $4$ or $7$. What is the largest number of elements $S^{}_{}$ can have?


 
 \textbf{Answer: }905



\textbf{Solution}

We first show that we can choose at most 5 numbers from $\{1, 2, \ldots , 11\}$ such that no two numbers have a difference of $4$ or $7$. We take the smallest number to be $1$, which rules out $5,8$. Now we can take at most one from each of the pairs: $[2,9]$, $[3,7]$, $[4,11]$, $[6,10]$. Now, $1989 = 180\cdot 11 + 9$. Because this isn't an exact multiple of $11$, we need to consider some numbers separately. 

 Notice that $1969 = 180\cdot11 - 11 = 179\cdot11$ (*). Therefore we can put the last $1969$ numbers into groups of 11. Now let's examine $\{1, 2, \ldots , 20\}$. If we pick $1, 3, 4, 6, 9$ from the first $11$ numbers, then we're allowed to pick $11 + 1$, $11 + 3$, $11 + 4$, $11 + 6$, $11 + 9$. This means we get 10 members from the 20 numbers. Our answer is thus $179\cdot 5 + 10 = \boxed{905}$.

 Remarks(*) Suppose that you choose the values 1, 2, 4, 7, and 10. Because $1989 = 180 \times 11 + 9$, this is not maximized. It is only maximized if we include the last element of the final set of 11, which is 10 (this is $\text{mod}(11)$ btw). To include the final element, we have to rewrite 1989 as $179 \times 11 + 20$, which includes the final element and increases our set by 1 element.

 




	\item (\textit{AIME I 1989}) Given a positive integer $n^{}_{}$, it can be shown that every complex number of the form $r+si^{}_{}$, where $r^{}_{}$ and $s^{}_{}$ are integers, can be uniquely expressed in the base $-n+i^{}_{}$ using the integers $0,1,2^{}_{},\ldots,n^2$ as digits. That is, the equation


 $r+si=a_m(-n+i)^m+a_{m-1}(-n+i)^{m-1}+\cdots +a_1(-n+i)+a_0$is true for a unique choice of non-negative integer $m^{}_{}$ and digits $a_0,a_1^{},\ldots,a_m$ chosen from the set $\{0^{}_{},1,2,\ldots,n^2\}$, with $a_m\ne 0^{}{}$. We write


 $r+si=(a_ma_{m-1}\ldots a_1a_0)_{-n+i}$to denote the base $-n+i^{}_{}$ expansion of $r+si^{}_{}$. There are only finitely many integers $k+0i^{}_{}$ that have four-digit expansions


 $k=(a_3a_2a_1a_0)_{-3+i^{}_{}}~~~~a_3\ne 0.$Find the sum of all such $k^{}_{}$.


 
 \textbf{Answer: }490



\textbf{Solution}

First, we find the first three powers of $-3+i$:

 $(-3+i)^1=-3+i ; (-3+i)^2=8-6i ; (-3+i)^3=-18+26i$

 So we solve the \href{https://artofproblemsolving.com/wiki/index.php?title=Diophantine\_equation}{diophantine equation} $a_1-6a_2+26a_3=0 \Longrightarrow a_1-6a_2=-26a_3$.

 The minimum the left-hand side can go is -54, so $1\leq a_3 \leq 2$ since $a_3$ can't equal 0, so we try cases:

 
                \begin{itemize}
                    \item Case 1: $a_3=2$
                \end{itemize}


                
            \begin{enumerate}[label=(\roman*)]
                \item The only solution to that is $(a_1, a_2, a_3)=(2,9,2)$.
            \end{enumerate}
            
                \begin{itemize}
                    \item Case 2: $a_3=1$
                \end{itemize}


                
            \begin{enumerate}[label=(\roman*)]
                \item The only solution to that is $(a_1, a_2, a_3)=(4,5,1)$.
            \end{enumerate}
            So we have four-digit integers $(292a_0)_{-3+i}$ and $(154a_0)_{-3+i}$, and we need to find the sum of all integers $k$ that can be expressed by one of those.

 $(292a_0)_{-3+i}$:

 We plug the first three digits into base 10 to get $30+a_0$. The sum of the integers $k$ in that form is $345$.

 $(154a_0)_{-3+i}$:

 We plug the first three digits into base 10 to get $10+a_0$. The sum of the integers $k$ in that form is $145$. The answer is $345+145=\boxed{490}$.~minor edit by \href{https://artofproblemsolving.com/wiki/index.php?title=User:Yiyj1&action=edit&redlink=1}{Yiyj1}

 




	\item (\textit{AIME I 1989}) Point $P$ is inside $\triangle ABC$. Line segments $APD$, $BPE$, and $CPF$ are drawn with $D$ on $BC$, $E$ on $AC$, and $F$ on $AB$ (see the figure below). Given that $AP=6$, $BP=9$, $PD=6$, $PE=3$, and $CF=20$, find the area of $\triangle ABC$.


 
                    \begin{figure}[H]
                        \centering
                        \includegraphics[width=0.5\textwidth]{Problem/AIME_I_1989_15_Problem_Figure_1.png}
                    \end{figure}
                    
 \textbf{Answer: }108



\textbf{Solution 1 (Ceva's Theorem, Stewart's Theorem)}

Let $[RST]$ be the area of polygon $RST$.  We'll make use of the following fact: if $P$ is a point in the interior of triangle $XYZ$, and line $XP$ intersects line $YZ$ at point $L$, then $\dfrac{XP}{PL} = \frac{[XPY] + [ZPX]}{[YPZ]}.$

 
                    \begin{figure}[H]
                        \centering
                        \includegraphics[width=0.5\textwidth]{Solution/AIME_I_1989_15_Solution_Figure_1.png}
                    \end{figure}
                    This is true because triangles $XPY$ and $YPL$ have their areas in ratio $XP:PL$ (as they share a common height from $Y$), and the same is true of triangles $ZPY$ and $LPZ$.

 We'll also use the related fact that $\dfrac{[XPY]}{[ZPX]} = \dfrac{YL}{LZ}$.  This is slightly more well known, as it is used in the standard proof of \href{https://artofproblemsolving.com/wiki/index.php?title=Ceva\%27s\_theorem}{Ceva's theorem}.

 Now we'll apply these results to the problem at hand.

 
                    \begin{figure}[H]
                        \centering
                        \includegraphics[width=0.5\textwidth]{Solution/AIME_I_1989_15_Solution_Figure_2.png}
                    \end{figure}
                    Since $AP = PD = 6$, this means that $[APB] + [APC] = [BPC]$; thus $\triangle BPC$ has half the area of $\triangle ABC$.  And since $PE = 3 = \dfrac{1}{3}BP$, we can conclude that $\triangle APC$ has one third of the combined areas of triangle $BPC$ and $APB$, and thus $\dfrac{1}{4}$ of the area of $\triangle ABC$.  This means that $\triangle APB$ is left with $\dfrac{1}{4}$ of the area of triangle $ABC$:
 \[[BPC]: [APC]: [APB] = 2:1:1.\]Since $[APC] = [APB]$, and since $\dfrac{[APC]}{[APB]} = \dfrac{CD}{DB}$, this means that $D$ is the midpoint of $BC$.

 Furthermore, we know that $\dfrac{CP}{PF} = \dfrac{[APC] + [BPC]}{[APB]} = 3$, so $CP = \dfrac{3}{4} \cdot CF = 15$.

 We now apply \href{https://artofproblemsolving.com/wiki/index.php?title=Stewart\%27s\_theorem}{Stewart's theorem} to segment $PD$ in $\triangle BPC$—or rather, the simplified version for a median.  This tells us that
 \[2 BD^2 + 2 PD^2 = BP^2+ CP^2.\] Plugging in we know, we learn that
 \begin{align*} 2 BD^2 + 2 \cdot 36 &= 81 + 225 = 306, \\ BD^2 &= 117. \end{align*}
Happily, $BP^2 + PD^2 = 81 + 36$ is also equal to 117.  Therefore $\triangle BPD$ is a right triangle with a right angle at $P$; its area is thus $\dfrac{1}{2} \cdot 9 \cdot 6 = 27$.  As $PD$ is a median of $\triangle BPC$, the area of $BPC$ is twice this, or 54.  And we already know that $\triangle BPC$ has half the area of $\triangle ABC$, which must therefore be $\boxed{108}$.

 





\textbf{Solution 2 [Quick] (Mass Points, Trig)}

We first assign masses:  Let $M(B) = 1$.  Then: $M(E) = 3$, $M(P) = 4$, $M(A) = M(D) = 2$, $M(C) = 1$, $M(F) = 3$.

 The masses on cevian $CF$ tell us that $CP = 15$.

 Also, $CD = DB$. Let $CD = DB = y$.

 Let angle $\angle PDB = \theta$.  Then, using Law of Cosines on triangle $PDB$, we find:
 \[36 + y^2 - 12y\cos(\theta) = 81\]

 Using Law of Cosines on triangle $CDP$, we find:
 \[36 + y^2 - 12\cos(180^\circ - \theta) = 225 \Rightarrow 36 + y^2 + 12\cos(\theta) = 225\]

 Therefore,
 \[36 + y^2 = \frac{225 + 81}{2} = 153 \Rightarrow y = \sqrt{117} = 3\sqrt{13}\]

 Next,
 \[12y\cos(\theta) = \frac{225 - 81}{2} = 72 \Rightarrow y\cos(\theta) = 6 \Rightarrow \cos(\theta) = \frac{2}{\sqrt{13}}\]

 Then,
 \[\cos(\theta) = \frac{2}{\sqrt{13}} \Rightarrow \sin(\theta) = \frac{3}{\sqrt{13}}\]

 Via the sine area formula:
 \[[PDB] = [PDC] = \frac{6y\sin(\theta)}{2} \Rightarrow [PDB] + [PDC] = [CDB] = 6y\sin(\theta)\]

 Then,
 \[[CAB] = 2 \cdot [CDB] = 12y\sin(\theta)\]

 Plugging in our values for $y$ and $\sin(\theta)$:
 \[12 \cdot 3\sqrt{13} \cdot \frac{3}{\sqrt{13}} = \boxed{108}\]

 





\textbf{Solution 3 (Mass Points, Stewart's Theorem, Heron's Formula)}

Because we're given three concurrent \href{https://artofproblemsolving.com/wiki/index.php?title=Cevian}{cevians} and their lengths, it seems very tempting to apply \href{https://artofproblemsolving.com/wiki/index.php?title=Mass\_points}{Mass points}. We immediately see that $w_E = 3$, $w_B = 1$, and $w_A = w_D = 2$. Now, we recall that the masses on the three sides of the triangle must be balanced out, so $w_C = 1$ and $w_F = 3$. Thus, $CP = 15$ and $PF = 5$. 

 Recalling that $w_C = w_B = 1$, we see that $DC = DB$ and $DP$ is a \href{https://artofproblemsolving.com/wiki/index.php?title=Median}{median} to $BC$ in $\triangle BCP$. Applying \href{https://artofproblemsolving.com/wiki/index.php?title=Stewart\%27s\_Theorem}{Stewart's Theorem}, we have the following:
 \[\frac{BC}{2}(9^2+15^2)=BC(6^2+ \left(\frac{BC}{2} \right)^2).\]Eliminating $BC$ on both sides, we have:
 \[\frac 12(9^2+15^2)=6^2+ \left(\frac{BC}{2} \right)^2.\]Combining terms and simplifying numbers, we have:
 \[153=36+\left(\frac{BC}{2} \right)^2.\]Subtracting 36 to the other side yields:
 \[117= \left(\frac{BC}{2} \right)^2.\]Finishing it off from there, we find that $BC=2 \sqrt{117}.$ Now, notice that $2[BCP] = [ABC]$, because both triangles share the same base, $BC$ and $h_{\triangle ABC} = 2h_{\triangle BCP}$. Applying \href{https://artofproblemsolving.com/wiki/index.php?title=Heron\%27s\_formula}{Heron's formula} on triangle $BCP$ with sides $15$, $9$, and $2\sqrt{117}$, we have:
 \[\sqrt{(\sqrt{117}+12)(\sqrt{117}+12-9)(\sqrt{117}+12-15)(\sqrt{117}+12-2\sqrt{117})}.\]Combining terms results in:
 \[\sqrt{(\sqrt{117}+12)(\sqrt{117}+3)(\sqrt{117}-3)(-\sqrt{117}+12)}.\]Notice that these factors can be grouped into a difference of squares:
 \[\sqrt{(144-\sqrt{117}^2)(\sqrt{117}^2-9)}.\]Since $\sqrt{117}^2=117$, we have:
 \[\sqrt{(27)(108)}.\]After simplifying this radical, we find that it equals $54.$ Therefore, $[BCP] = 54$, and hence $[ABC]=2 \cdot 54= \boxed{108}$.

 (The original author made a mistake in their solution. Corrected and further explained by dbnl.)

 





\textbf{Solution 4 (Ceva's Theorem, Stewart's Theorem)}

Using a different form of \href{https://artofproblemsolving.com/wiki/index.php?title=Ceva\%27s\_Theorem}{Ceva's Theorem}, we have $\frac {y}{x + y} + \frac {6}{6 + 6} + \frac {3}{3 + 9} = 1\Longleftrightarrow\frac {y}{x + y} = \frac {1}{4}$

 Solving $4y = x + y$ and $x + y = 20$, we obtain $x = CP = 15$ and $y = FP = 5$.

 Let $Q$ be the point on $AB$ such that $FC \parallel QD$.Since $AP = PD$ and $FP\parallel QD$, $QD = 2FP = 10$. (Stewart's Theorem)

 Also, since $FC\parallel QD$ and $QD = \frac{FC}{2}$, we see that $FQ = QB$, $BD = DC$, etc. (\href{https://artofproblemsolving.com/wiki/index.php?title=Stewart\%27s\_Theorem}{Stewart's Theorem})Similarly, we have $PR = RB$ ($= \frac12PB = 7.5$) and thus $RD = \frac12PC = 4.5$.

 $PDR$ is a $3-4-5$ \href{https://artofproblemsolving.com/wiki/index.php?title=Right\_triangle}{right triangle}, so $\angle PDR$ ($\angle ADQ$) is $90^\circ$.Therefore, the area of $\triangle ADQ = \frac12\cdot 12\cdot 6 = 36$.Using area ratio, $\triangle ABC = \triangle ADB\times 2 = \left(\triangle ADQ\times \frac32\right)\times 2 = 36\cdot 3 = \boxed{108}$.

 





\textbf{Solution 5 (Stewart's Theorem)}

First, let $[AEP]=a, [AFP]=b,$ and $[ECP]=c.$ Thus, we can easily find that $\frac{[AEP]}{[BPD]}=\frac{3}{9}=\frac{1}{3} \Leftrightarrow [BPD]=3[AEP]=3a.$ Now, $\frac{[ABP]}{[BPD]}=\frac{6}{6}=1\Leftrightarrow [ABP]=3a.$ In the same manner, we find that $[CPD]=a+c.$ Now, we can find that $\frac{[BPC]}{[PEC]}=\frac{9}{3}=3 \Leftrightarrow \frac{(3a)+(a+c)}{c}=3 \Leftrightarrow c=2a.$ We can now use this to find that $\frac{[APC]}{[AFP]}=\frac{[BPC]}{[BFP]}=\frac{PC}{FP} \Leftrightarrow \frac{3a}{b}=\frac{6a}{3a-b} \Leftrightarrow a=b.$ Plugging this value in, we find that $\frac{FC}{FP}=3 \Leftrightarrow PC=15, FP=5.$ Now, since $\frac{[AEP]}{[PEC]}=\frac{a}{2a}=\frac{1}{2},$ we can find that $2AE=EC.$ Setting $AC=b,$ we can apply \href{https://artofproblemsolving.com/wiki/index.php?title=Stewart\%27s\_Theorem}{Stewart's Theorem} on triangle $APC$ to find that $(15)(15)(\frac{b}{3})+(6)(6)(\frac{2b}{3})=(\frac{2b}{3})(\frac{b}{3})(b)+(b)(3)(3).$ Solving, we find that $b=\sqrt{405} \Leftrightarrow AE=\frac{b}{3}=\sqrt{45}.$ But, $3^2+6^2=45,$ meaning that $\angle{APE}=90 \Leftrightarrow [APE]=\frac{(6)(3)}{2}=9=a.$ Since $[ABC]=a+a+2a+2a+3a+3a=12a=(12)(9)=108,$ we conclude that the answer is $\boxed{108}$.

 





\textbf{Solution 6 (Mass of a Point, Stewart's Theorem, Heron's Formula)}

Firstly, since they all meet at one single point, denoting the mass of them separately. Assuming $M(A)=6;M(D)=6;M(B)=3;M(E)=9$; we can get that $M(P)=12;M(F)=9;M(C)=3$; which leads to the ratio between segments, 
 \[\frac{CE}{AE}=2;\frac{BF}{AF}=2;\frac{BD}{CD}=1.\] Denoting that $CE=2x;AE=x; AF=y; BF=2y; CD=z; DB=z.$ 

 Now we know three cevians' length, Applying Stewart theorem to them, getting three different equations:
 \begin{align} (3x)^2 \cdot 2y+(2z)^2 \cdot y&=(3y)(2y^2+400), \\ (3y)^2 \cdot z+(3x)^2 \cdot z&=(2z)(z^2+144), \\ (2z)^2 \cdot x+(3y)^2 \cdot x&=(3x)(2x^2+144). \end{align}After solving the system of equation, we get that $x=3\sqrt{5};y=\sqrt{13};z=3\sqrt{13}$;

 pulling $x,y,z$ back to get the length of $AC=9\sqrt{5};AB=3\sqrt{13};BC=6\sqrt{13}$; now we can apply Heron's formula here, which is 
 \[\sqrt\frac{(9\sqrt{5}+9\sqrt{13})(9\sqrt{13}-9\sqrt{5})(9\sqrt{5}+3\sqrt{13})(9\sqrt{5}-3\sqrt{13})}{16}=108.\]

 Our answer is $\boxed{108}$.

 \textbf{Note (how to find x and y without the system of equations)} To ease computation, we can apply Stewart's Theorem to find $x$, $y$, and $z$ directly. Since $M(C)=3$ and $M(F)=9$, $\overline{PC}=15$ and $\overline{PF}=5$. We can apply Stewart's Theorem on $\triangle CPE$ to get $(2x+x)(2x \cdot x) + 3^2 \cdot 3x = 15^2 \cdot x + 6^2\cdot 2x$. Solving, we find that $x=3\sqrt{5}$. We can do the same on $\triangle APB$ and $\triangle BPC$ to obtain $y$ and $z$. We proceed with Heron's Formula as the solution states. 

 





\textbf{Solution 7 (easier version of Solution 5)}

In Solution 5, instead of finding all of $x, y, z$, we only need $y, z$. This is because after we solve for $y, z$, we can notice that $\triangle BAD$ is isosceles with $AB = BD$. Because $P$ is the midpoint of the base, $BP$ is an altitude of $\triangle BAD$. Therefore, $[BAD] = \frac{(AD)(BP)}{2} = \frac{12 \cdot 9}{2} = 54$. Using the same altitude property, we can find that $[ABC] = 2[BAD] = 2 \cdot 54 = \boxed{108}$.

 -NL008

 





\textbf{Solution 8 (Mass Points, Stewart's Theorem, Simple Version)}

Set $AF=x,$ and use \href{https://artofproblemsolving.com/wiki/index.php?title=Mass\_points}{mass points} to find that $PF=5$ and $BF=2x.$ Using \href{https://artofproblemsolving.com/wiki/index.php?title=Stewart\%27s\_Theorem}{Stewart's Theorem} on $APB,$ we find that $AB=3\sqrt{13}.$ Then we notice that $APB$ is right, which means the area of $APB$ is $27.$ Because $CF=4\cdot PF,$ the area of $ABC$ is $4$ times the area of $APB,$ which means the area of $ABC=4\cdot 27=\boxed{108}.$

 





\textbf{Solution 9 (Ratios, Auxiliary Lines and 3-4-5 triangle)}

We try to solve this using only elementary concepts. Let the areas of triangles $BCP$, $ACP$ and $ABP$ be $X$, $Y$ and $Z$ respectively. Then $\frac{X}{Y+Z}=\frac{6}{6}=1$ and $\frac{Y}{X+Z}=\frac{3}{9}=\frac{1}{3}$. Hence $\frac{X}{2}=Y=Z$. Similarly $\frac{FP}{PC}=\frac{Z}{X+Y}=\frac{1}{3}$ and since $CF=20$ we then have $FP=5$. Additionally we now see that triangles $FPE$ and $CPB$ are similar, so $FE \parallel BC$ and $\frac{FE}{BC} = \frac{1}{3}$. Hence $\frac{AF}{FB}=\frac{1}{2}$. Now construct a point $K$ on segment $BP$ such that $BK=6$ and $KP=3$, we will have $FK \parallel AP$, and hence $\frac{FK}{AP} = \frac{FK}{6} = \frac{2}{3}$, giving $FK=4$. Triangle $FKP$ is therefore a 3-4-5 triangle! So $FK \perp BE$ and so $AP \perp BE$. Then it is easy to calculate that $Z = \frac{1}{2} \times 6 \times 9 = 27$ and the area of triangle $ABC = X+Y+Z = 4Z = 4 \times 27 = \boxed{108}$.~Leole

 


 





\textbf{Solution 10 (Just Trig Bash)}

We start with mass points as in Solution 2, and receive $BF:AF = 2$, $BD:CD = 1$, $CE:AE = 2$. \href{https://artofproblemsolving.com/wiki/index.php?title=Law\_of\_Cosines}{Law of Cosines} on triangles $ADB$ and $ADC$ with $\theta = \angle ADB$ and $BD=DC=x$ gives
 \[36+x^2-12x\cos \theta = 81\]
 \[36+x^2-12x\cos (180-\theta) = 36+x^2+12x\cos \theta = 225\]Adding them: $72+2x^2=306 \implies x=3\sqrt{13}$, so $BC = 6\sqrt{13}$. Similarly, $AB = 3\sqrt{13}$ and $AC = 9\sqrt{5}$. Using Heron's,
 \[[ \triangle ABC ]= \sqrt{\left(\dfrac{9\sqrt{13}+9\sqrt{5}}{2}\right)\left(\dfrac{9\sqrt{13}+9\sqrt{5}}{2}\right)\left(\dfrac{3\sqrt{13}+9\sqrt{5}}{2}\right)\left(\dfrac{-3\sqrt{13}+9\sqrt{5}}{2}\right)} = \boxed{108}.\]

 





\textbf{Solution 11 (Assumption)}

APB looks really close to a right triangle so we use the Pythagorean Theorem and we get $AB=\sqrt{117}$; therefore the area of the triangle is equal to $\frac{\sqrt{117}\cdot 20}{2}$, which is around $\boxed{108}$.

 


 




\end{enumerate}